\documentclass[letterpaper,10pt]{article}
\usepackage[left=2cm,top=3cm,right=2cm,bottom=4cm]{geometry} 
\usepackage[utf8]{inputenc}
\usepackage[activeacute,spanish]{babel}
\usepackage{multirow}	% Para poder unir filas en las tablas
\usepackage{colortbl}	% Para colorear tablas
\usepackage{latexsym,amssymb,oldlfont,amsmath}
\usepackage{graphicx,graphics}
\usepackage{multicol}
\usepackage{color}
\usepackage{listings}
\usepackage{lastpage}
\usepackage{enumitem}
\usepackage{mathtools}


%\usepackage{fontspec}
%\setmainfont{Times New Roman}
\newcommand{\R}{{\mathbb R}}
\newcommand{\si}{{\mathcal S}^n}
\renewcommand{\theequation}{\thesection.\arabic{equation}} % Cambia formato salida numeracion equation
\makeatletter\@addtoreset {equation}{section} \makeatother % Tambien se usa
\renewcommand{\baselinestretch}{1.3}
\usepackage{amssymb}
\usepackage{hyperref}
\definecolor{midnightblue}{rgb}{0.1,0.1,0.44}


\newcommand{\dis}{\displaystyle}
\renewcommand{\theenumii}{\arabic{enumii}}

\def\Z{\mathbb{Z}}
\def\N{\mathbb{N}}
\def\R{\mathbb{R}}
\def\M{\mathcal{M}_n(\mathbb{R})}
\def\sin{\operatorname{sen}}
\def\e{\operatorname{e}}
\DeclareMathOperator{\arcsec}{arcsec}
\DeclareMathOperator{\arccot}{arccot}
\DeclareMathOperator{\arccsc}{arccsc}
\providecommand{\abs}[1]{\lvert#1\rvert}


\usepackage{etoolbox,fancyhdr,xcolor}
\fancyhf{} % sets both header and footer to nothing
\renewcommand{\headrulewidth}{0pt}
%\renewcommand{\footrulewidth}{0.5pt}
\pagestyle{fancy}
\headheight=60pt %para cambiar el tamaño del encabezadox
\fancyhead[L] %la "L" indica a la izquierda
{
\begin{minipage}{6cm}
\includegraphics[width=0.9\textwidth]{logo_udp.eps}
\end{minipage}
}


\fancyfoot[C]{}

\title{C2_Ay7_pauta}

\begin{document}

\begin{center}
Cálculo II\\
Ayudantía Nº7 - Pauta\\
Primer Semestre 2017
\end{center}

\begin{enumerate}
\item Determine si las siguientes afirmaciones son verdaderas o falsas y justifique su respuesta:

\begin{enumerate}[label=\emph{\alph*)}]

\item El valor de $\displaystyle \int_0^\pi f(x) \; dx$, para $f(x)= \left\{ \begin{array}{lcc}
             \sin x &   si  & 0 \leq x < \frac{\pi}{2} \\
             \\ \cos x &  si & \frac{\pi}{2} \leq x \leq \pi 
             \end{array}
   \right.$ es $2\pi$\\


\begin{align*}
\dis \int_0^\pi f(x)dx &= \int_0^\frac{\pi}{2} \sin (x)dx +\int_\frac{\pi}{2}^\pi \cos(x)dx\\
\dis  &=  -\displaystyle  \left.\cos(x)\right|_{0}^{\frac{\pi}{2}} + \displaystyle  \left.\sin(x)\right|_{\frac{\pi}{2}}^{\pi}\\
\dis  &= -(0-1) + (0-1)=0\\ 
\dis 0 &\neq	2\pi
\end{align*}

Por lo tanto la afirmación en \underline{FALSA}.

\dotfill\textbf{[1 punto]}

\item El área limitada por las curvas $y = \frac{x}{\sqrt{2x^2+1}}$; $x=0$; $x=2$; $y=0$ es $1$.\\

Debemos calcular la siguiente integral $\dis \int_0^2 \frac{x}{\sqrt{2x^2+1}}dx$.\\

Para esto aplicamos la sustitución $u=2x^2+1$ y $du=4xdx$ obteniendo lo siguiente:

\begin{align*}
\dis \int_0^2 \frac{x}{\sqrt{2x^2+1}}dx &= \frac{1}{4}\int_1^9 \frac{1}{\sqrt{u}}du\\
\dis  &= \displaystyle  \frac{2}{4} \left.\sqrt{u}\right|_{1}^{9}\\
\dis  &= \frac{3}{2}- \frac{1}{2} = 1\\
\end{align*}
Por lo tanto la afirmación en \underline{VERDADERA}.

\dotfill\textbf{[1 punto]}

\newpage
\item El valor de $c \in R$, tal que el área acotada por $y = x^2+c$; $x=0$; $x=3$; $y=0$ sea 12, es $c=1$.\\

\begin{align*}
\dis \int_0^3 (x^2 +1)dx &= \displaystyle   \left.(\frac{x^3}{3}+x)\right|_{0}^{3}\\
\dis  &= 9+3\\
\dis  &= 12
\end{align*}

Por lo tanto la afirmación en \underline{VERDADERA}.

\dotfill\textbf{[1 punto]}

\newpage

\end{enumerate}

\item Grafique la regi\'{o}n  limitada por las curvas y calcule su \'{a}rea\\
$\dis y=\frac{1}{x}$ \, ;\,  $\dis y=\frac{1}{x^2}$ \, ;\, $x=1$ \, ;\, $x=2$

\begin{figure}[!htb]\centering
   \begin{minipage}{0.6\textwidth}
     \frame{\includegraphics[width=8cm]{c2-ay7-1.png}}
     
   \end{minipage}
   \begin {minipage}{0.3\textwidth}
     \frame{\includegraphics[width=2cm]{c2-ay7-11.png}}
     
   \end{minipage}
\end{figure}

\begin{align*}
\dis \int_1^2 (\frac{1}{x}- \frac{1}{x^2})dx &= \displaystyle   \left.(\ln (x)+\frac{1}{x})\right|_{1}^{2}\\
\dis  &= \ln 2 + \frac{1}{2} - (\ln 1 + 1)\\
\dis  &= \ln 2 - \frac{1}{2}
\end{align*}

\dotfill\textbf{[1 punto]}

\item Grafique la regi\'{o}n  limitada por las curvas y calcule su \'{a}rea\\
$\dis y=\e^x$ \, ;\,  $\dis y=\e^{-x}$ \, ;\, $x=-2$ \, ;\, $x=1$ \, ;\,  $\dis y=0$


\begin{figure}[!htb]\centering
   \begin{minipage}{0.6\textwidth}
     \frame{\includegraphics[width=8cm]{c2-ay7-2.png}}
     
   \end{minipage}
   \begin {minipage}{0.3\textwidth}
     \frame{\includegraphics[width=2cm]{c2-ay7-21.png}}
     
   \end{minipage}
\end{figure}

\begin{align*}
\dis \int_{-2}^0 e^{x}dx + \int_{0}^1 e^{-x}dx &= \displaystyle   \left.e^{x}\right|_{-2}^{0} - \left.e^{-x}\right|_{0}^{1} = 1 - e^{-2}-e^{-1}+1 \\
\dis  &= 2 - e^{-2}-e^{-1}
\end{align*}

\dotfill\textbf{[1 punto]}


\newpage
\item Calcule el \'{a}rea de la regi\'{o}n limitada por la curva
$y=x^4-2x^3+x^2+3$, el eje $X$ y las rectas $x=a$ ; $x=b$ donde  $a$
y $b$ son las abscisas de los puntos m\'{\i}nimos de la funci\'{o}n.\\

Calculamos los mínimos de la función  $y=x^4-2x^3+x^2+3$:\\
\begin{align*}
\dis y &= x^4-2x^3+x^2+3\\
\dis y' &= 4x^3- 6x^2+2x =0\\
\dis y' &= 2x(x-1)(2x-1)
\end{align*}
Los puntos críticos son $\{0,\frac{1}{2},1\}$. Ahora,
\begin{align*}
\dis f''(x) &=12x^2-12x+2\\
\dis f''(0) &=12x^2-12x+2 = 2\\
\dis  f''(\frac{1}{2}) &=12x^2-12x+2 = -1\\
\dis  f''(1) &=12x^2-12x+2 = 2\\
\end{align*}

Por lo tanto los mínimos ocurren en $x =\{0,1\}$.

\dotfill\textbf{[1 punto]}

Luego calculamos la región solicitada.

\begin{align*}
\dis \int_{0}^1 (x^4-2x^3+x^2+3)dx &= \displaystyle   \left.(\frac{x^5}{5}- \frac{x^4}{2} + \frac{x^3}{3} + 3x)\right|_{0}^{1}\\
\dis  &= \frac{1}{5} - \frac{1}{2} + \frac{1}{3} + 3\\
\dis  &= \frac{91}{30}
\end{align*}

\dotfill\textbf{[1 punto]}



\newpage
\item Calcule el \'{a}rea de la regi\'{o}n limitada por la par\'{a}bola
$y=x^2-2x+2$, su tangente en el punto $(3,5)$ y el eje $Y$.\\


Primero calculamos la pendiente que toma la recta en el punto $(3,5)$ derivando la función $y=x^2-2x+2$.

\begin{align*}
\dis f(x) &= x^2-2x+2\\
\dis f'(x) &= 2x-2\\
\dis f'(3) &= 4\\
\end{align*}

Formamos la recta con la pendiente obtenida y el punto $(3,5)$.

\begin{align*}
\dis (y-5) &= 4(x-3)\\
\dis  y &= 4x-7
\end{align*}

\dotfill\textbf{[1 punto]}

Como referencia graficamos las funciones:

\begin{figure}[!htb]\centering
   \begin{minipage}{0.6\textwidth}
     \frame{\includegraphics[width=5cm]{c2-ay7-3.png}}
     
   \end{minipage}
   \begin {minipage}{0.3\textwidth}
     \frame{\includegraphics[width=2cm]{c2-ay7-31.png}}
     
   \end{minipage}
\end{figure}


\newpage
Finalmente calculamos la integral

\begin{align*}
\dis \int_{0}^3 ((x^2-2x+2) - (4x-7))dx &= \int_{0}^3 (x^2-6x+9)dx\\
\dis  &= \displaystyle   \left.(\frac{x^3}{3} + \frac{6x^2}{2} + 9x)\right|_{0}^{3}\\
\dis  &= 9-27+27 = 9
\end{align*}

\dotfill\textbf{[1 punto]}


\end{enumerate}




\vspace{10mm}




\end{document}
